\documentclass[11pt,a4paper]{article}
\usepackage[utf8]{inputenc}
\usepackage[T1]{fontenc}
\usepackage[portuguese]{babel}
\usepackage[a4paper,margin=2.5cm]{geometry}
\usepackage{booktabs}
\usepackage{tabularx}
\usepackage{hyperref}
\usepackage{xurl}
\usepackage{enumitem}
\usepackage{verbatim}

\hypersetup{colorlinks=true,linkcolor=black,urlcolor=blue}
\setlist{noitemsep,topsep=3pt}

\title{Segurança de Dados\\\large Assignment 2 -- Classical Cryptographic Schemes}
\author{Hugo Soares}
\date{8 de outubro de 2026}

\begin{document}
\maketitle

\section*{Questão 1}

Foi implementado um ataque a uma cifra monoalfabética utilizando o texto cifrado
\texttt{Q1\_ciphertext.txt}. O texto original está escrito em português e todos os
caracteres que não pertencem ao alfabeto latino foram removidos, conforme indicado
no enunciado.

\subsection*{Método utilizado}

Foi seguido um procedimento manual, apoiado por um programa simples. O
programa \texttt{monoalphabetic\_attack.py} não procura automaticamente a
chave: apenas conta as letras do texto, aplica as correspondências que foram
deduzidas e mostra o resultado parcial.

O primeiro passo foi remover os caracteres que não pertencem ao alfabeto
\texttt{A--Z} e contar as ocorrências de cada letra. A percentagem de uma letra
\texttt{x} foi calculada por:

\[
\mathrm{freq}(x) =
\frac{\mathrm{ocorrencias}(x)}
{\mathrm{numero\ total\ de\ letras}}\times 100.
\]

As letras mais frequentes encontradas no texto cifrado foram:

\begin{center}
\begin{tabular}{c|c|c}
\textbf{Letra cifrada} & \textbf{Ocorrências} & \textbf{Percentagem}\\
\midrule
\texttt{I} & 322 & 14,25\%\\
\texttt{J} & 298 & 13,19\%\\
\texttt{K} & 249 & 11,02\%\\
\texttt{T} & 172 & 7,61\%\\
\texttt{Z} & 151 & 6,68\%\\
\end{tabular}
\end{center}

Estas frequências foram comparadas com as frequências esperadas das letras em
português. Como a cifra é monoalfabética, a letra mais frequente do texto
cifrado é uma candidata à letra mais frequente do português. Esta comparação
não foi considerada suficiente por si só, porque frequências próximas podem
corresponder a letras diferentes num texto específico.

As hipóteses foram, por isso, verificadas através de digramas, trigramas e do
contexto. Por exemplo, o início do texto cifrado é:

\begin{verbatim}
JTJZHJTIKTCJZKI...
\end{verbatim}

Ao testar que \texttt{JTJ} correspondia a \texttt{ASA}, obtive:

\begin{verbatim}
JTJZHJTIKT
ASARCASEOS
\end{verbatim}

Esta hipótese produz uma sequência compatível com o início de
\textit{Os Lusíadas}, ``AS ARMAS E OS...''. A partir desta observação foram
registadas as correspondências \texttt{J--A} e \texttt{T--S}. As letras
seguintes foram deduzidas de forma semelhante e confirmadas sempre que
apareciam noutras posições do texto.

O programa executa apenas os seguintes passos:
\begin{enumerate}
\item lê o ficheiro, converte o conteúdo para maiúsculas e mantém apenas letras
      entre \texttt{A} e \texttt{Z};
\item mostra a contagem e a percentagem de cada letra;
\item aplica a tabela \texttt{CHAVE}, que foi preenchida manualmente; e
\item grava o texto claro em \texttt{Q1\_plaintext.txt}, quando é indicado um
      segundo argumento.
\end{enumerate}

As tabelas de digramas e trigramas fornecidas no enunciado foram usadas como
referência durante esta análise. A frequência de letras serviu para propor
correspondências iniciais; os digramas e trigramas serviram para testar se essas
correspondências eram plausíveis. Uma correspondência só foi mantida quando
produzia sequências coerentes e respeitava a regra de que letras cifradas
diferentes devem corresponder a letras claras diferentes.

\subsection*{Chave e permutação descobertas}

Comparando o texto cifrado com o texto claro obtido, a permutação usada na
decifragem é:

\begin{center}
\small
\begin{tabularx}{\textwidth}{>{\ttfamily}c>{\ttfamily}c@{\quad} >{\ttfamily}c>{\ttfamily}c@{\quad} >{\ttfamily}c>{\ttfamily}c@{\quad} >{\ttfamily}c>{\ttfamily}c}
\toprule
\textbf{Cifrado} & \textbf{Claro} & \textbf{Cifrado} & \textbf{Claro} &
\textbf{Cifrado} & \textbf{Claro} & \textbf{Cifrado} & \textbf{Claro} \\
\midrule
A & J & B & C & C & B & D & P \\
E & I & F & Q & G & L & H & M \\
I & E & J & A & K & O & L & G \\
M & H & P & D & Q & F & R & X \\
S & T & T & S & U & V & V & U \\
W & Z & Y & N & Z & R & & \\
\bottomrule
\end{tabularx}
\end{center}

As letras \texttt{n}, \texttt{o} e \texttt{x} não aparecem no texto cifrado, pelo
que a sua correspondência não pode ser inferida a partir desta instância. Para
as letras presentes, o mapeamento é consistente e bijetivo. O texto contém
2259 letras, tanto no ficheiro cifrado como no ficheiro claro.

\subsection*{Resultado}

O início do texto claro recuperado é:

\begin{verbatim}
AS ARMAS E OS BAROES ASSINALADOS QUE DA OCIDENTAL PRAIA LUSITANA
\end{verbatim}

Depois de reintroduzir os espaços apenas para leitura, este trecho corresponde
ao início de \textit{Os Lusíadas}, de Luís de Camões. O resultado completo é
mantido no ficheiro \texttt{Q1\_plaintext.txt}; como o enunciado pede que os
espaços e caracteres especiais sejam removidos, o ficheiro contém a versão
contínua em maiúsculas.

O programa pode ser executado a partir da pasta \texttt{SD} com:

\begin{verbatim}
python3 Aula2/monoalphabetic_attack.py \
    Aula2/Q1_ciphertext.txt Aula2/Q1_plaintext.txt
\end{verbatim}

\section*{Questão 2}

O enunciado pede a implementação de um ataque a uma cifra de Vigenère,
recebendo o tamanho da chave, o caminho do texto cifrado e uma lista de palavras
que possam ocorrer no texto claro. Também pede a utilização de uma chave de
tamanho quatro sobre o ficheiro \texttt{Q2\_ciphertext.txt}.

Na versão fornecida da pasta da aula, esse ficheiro não está presente e também
não existe um programa \texttt{vigenere\_attack.py}. Assim, não é possível
produzir de forma verificável a chave, o texto claro ou os resultados
experimentais da Questão 2. Para completar esta parte, seria necessário
adicionar o ficheiro de entrada e implementar o ataque descrito no enunciado.

\subsection*{Ataque esperado}

Para cada uma das quatro posições da chave, devem ser extraídos os caracteres
do texto cifrado nessa posição. Cada sequência é uma cifra de César
independente. A hipótese de deslocamento de cada posição pode ser testada
decifrando as sequências e verificando se o texto reconstruído contém pelo
menos uma das palavras \texttt{DIREITO}, \texttt{FANTASIA},
\texttt{LIBERDADE}, \texttt{MUNDO}, \texttt{QUE} ou
\texttt{INTERVENCAO}. Quando uma combinação satisfizer essa condição, o
programa deve imprimir a chave na primeira linha e o texto claro completo na
segunda linha.

\section*{Conclusão}

A Questão 1 foi resolvida manualmente. A frequência das letras permitiu
estabelecer as primeiras hipóteses através da comparação com a frequência
esperada em português. Essas hipóteses foram confirmadas ou rejeitadas através
de digramas, trigramas e do contexto das sequências obtidas. O programa foi
mantido simples para tornar explícito o raciocínio: a chave está escrita no
código e o programa apenas aplica a substituição escolhida. O resultado obtido
é coerente com o início de \textit{Os Lusíadas}. A Questão 2 fica pendente por
falta dos ficheiros necessários na pasta fornecida.

\end{document}
